% from ega4-10-fr.tex
\oldpage[IV-3]{101}
\subsection{Espaces de Jacobson}
\label{subsection:IV.10.3-fr}

\begin{definition}[10.3.1]
\label{IV.10.3.1-fr}
On dit qu'un espace topologique $X$ est un \emph{espace de Jacobson} si
l'ensemble $X_0$ des points fermés de $X$ est très dense dans $X$ (autrement
dit, si l'injection canonique $X_0\to X$ est un
quasi-homéomorphisme).
\end{definition}

Cela signifie donc (10.1.2) que toute partie localement fermée (ou seulement
localement quasi-constructible) $Z\neq\varnothing$ de $X$ contient un point
fermé de $X$, ou encore que \emph{toute partie fermée de $X$ est l'adhérence
de l'ensemble de ses points fermés}.

\begin{proposition}[10.3.2]
\label{IV.10.3.2-fr}
Soient $X$ un espace de Jacobson, $Z$ une partie localement
quasi-constructible de $X$~; alors le sous-espace $Z$ de $X$ est un espace de
Jacobson, et pour qu'un point de $Z$ soit fermé dans $Z$, il faut et il suffit
qu'il soit fermé dans $X$.
\end{proposition}

Si $X_0$ est l'ensemble des points fermés de $X$, $Z\cap X_0$ est très dense
dans $Z$ en vertu de (10.2.4) appliqué à l'injection $i:X_0\to X$~; comme
l'ensemble $Z_0$ des points fermés de $Z$ contient évidemment $Z\cap X_0$,
il est \emph{a fortiori} très dense dans $Z$, donc $Z$ est un espace de
Jacobson. Montrons maintenant qu'en fait on a $Z_0=Z\cap X_0$~; soit $x\in Z$
un point fermé dans $Z$~; soit $\overline{\{x\}}$ son adhérence dans $X$~;
$\overline{\{x\}}\cap Z=\{x\}$ est donc une partie localement
quasi-constructible de $X$, et comme son intersection avec $X_0$ est non vide
(10.1.2), on a $x\in X_0$.

\begin{proposition}[10.3.3]
\label{IV.10.3.3-fr}
Soient $X$ un espace topologique, $(U_\alpha)$ un recouvrement ouvert de $X$.
Pour que $X$ soit un espace de Jacobson, il faut et il suffit que chacun des
sous-espaces $U_\alpha$ le soit.
\end{proposition}

La condition est nécessaire en vertu de (10.3.2). Inversement, montrons
d'abord que l'hypothèse que les $U_\alpha$ sont des espaces de Jacobson
entraîne que pour qu'un point $x\in U_\alpha$ soit fermé dans $X$, il suffit
qu'il soit fermé dans $U_\alpha$. Il suffit en effet de voir que cette
condition entraîne que $x$ est aussi fermé dans chacun des $U_\beta$ qui le
contiennent~; mais $U_\alpha\cap U_\beta$ est ouvert dans $U_\alpha$, donc $x$
est fermé dans $U_\alpha\cap U_\beta$, et en vertu de (10.3.2), $x$ est aussi
fermé dans $U_\beta$, ce qui achève la démonstration.

\subsection{Préschémas de Jacobson et anneaux de Jacobson}
\EGAReaderTarget{ega4.c04.subsec.10-4}{subsection:IV.10.4-fr}

\begin{definition}[10.4.1]
\label{IV.10.4.1-fr}
On dit qu'un préschéma $X$ est un \emph{préschéma de Jacobson} si l'espace
topologique $X$ sous-jacent est un espace de Jacobson. On dit qu'un anneau
$A$ est un \emph{anneau de Jacobson} si $\operatorname{Spec}(A)$ est un
préschéma de Jacobson.
\end{definition}

Toute partie fermée de $X=\operatorname{Spec}(A)$ est de la forme
$Z=V(\mathfrak a)$, où $\mathfrak a$ est un idéal égal à sa racine, et
l'ensemble $X_0$ des points fermés de $X$ est l'ensemble des idéaux maximaux
de $A$~; dire que $Z\cap X_0$ est dense dans $Z$ signifie donc que
$\mathfrak a$ est \emph{intersection d'idéaux maximaux} (\textbf{I}, 1.1.4)~;
comme $\mathfrak a$ est intersection d'idéaux premiers, il revient au même de
dire que tout idéal premier de $A$ est intersection d'idéaux maximaux~; en
vertu de (10.3.1) et (10.1.2), la définition usuelle des anneaux de Jacobson
(Bourbaki, \emph{Alg. comm.}, chap.~V, \S~3, n\textsuperscript{o}~4, déf.~1)
coïncide donc avec la définition (10.4.1).

